\section*{Hoofdstuk \ref{ch:g4:weight-and-mass} --- Gewicht en massa}

\begin{solution}{exo:g4:weight-and-mass:1}
\omterm{def:g3:measuring-mass:mass}{Massa}: hoeveel stof? \omterm{def:g4:weight-and-mass:weight}{Gewicht}: hoe hard wordt het hier omlaag
getrokken? \omterm{def:g4:weight-and-mass:weight}{Gewicht} is de \omterm{def:g2:pushes-and-pulls:force}{kracht} --- de neerwaartse aantrekking van
de Aarde.
\end{solution}

\begin{solution}{exo:g4:weight-and-mass:2}
Recht omlaag, altijd. De \omterm{def:g4:weight-and-mass:springscale}{veerunster} --- \omterm{def:g4:weight-and-mass:weight}{gewicht} rekt haar veer uit,
en de wijzer meldt de rek.
\end{solution}

\begin{solution}{exo:g4:weight-and-mass:3}
Ongeveer vier vingerbreedtes --- dubbel zoveel stof, dubbel zoveel
trek, dubbel zoveel rek: de propositie ``meer stof, sterkere trek''.
\end{solution}

\begin{solution}{exo:g4:weight-and-mass:4}
De rek moet \emph{vanaf de lege stand} gemeten worden: het elastiekje
hangt al een stukje uit onder het eigen \omterm{def:g4:weight-and-mass:weight}{gewicht} van het bekertje, en
alleen het extra zakken onder het nulstreepje meet de lading.
\end{solution}

\begin{solution}{exo:g4:weight-and-mass:5}
Ja --- \omterm{def:g3:measuring-mass:mass}{massa} reist onveranderd mee: nog steeds \qty{20}{kg} stof. Maar
hij is veel makkelijker omhoog te houden: de \omterm{def:g2:moon-in-the-sky:moon}{Maan} trekt er ongeveer
zes keer zwakker aan, dus is zijn \omterm{def:g4:weight-and-mass:weight}{gewicht} ongeveer een zesde van de
waarde op Aarde.
\end{solution}

\begin{solution}{exo:g4:weight-and-mass:6}
Hun benen duwen even hard als op Aarde getrainde benen altijd doen,
maar de tegentrek van de \omterm{def:g2:moon-in-the-sky:moon}{Maan} is zes keer zwakker --- dus wordt elke
gewone stap een sprong.
\end{solution}

\begin{solution}{exo:g4:weight-and-mass:7}
De \omterm{def:g4:weight-and-mass:springscale}{veerunster}: die voelt de plaatselijke trek en leest zes keer
minder af. De \omterm{def:g3:levers-and-scales:scale}{balansweegschaal} vergelijkt de trek op de rugzak met de
trek op de blokjes, en de \omterm{def:g2:moon-in-the-sky:moon}{Maan} verzwakt allebei even sterk --- de
schalen spelen nog steeds gelijk bij dezelfde \omterm{def:g3:measuring-mass:mass}{massa}.
\end{solution}

\begin{solution}{exo:g4:weight-and-mass:8}
Waar. \omterm{def:g4:weight-and-mass:weight}{Gewicht} is de aantrekking van de Aarde op de stof van een ding,
en die trek groeit met de hoeveelheid stof --- geen stof, niets om aan
te trekken: overal \omterm{def:g4:weight-and-mass:weight}{gewicht} nul.
\end{solution}

\begin{solution}{exo:g4:weight-and-mass:9}
De veer binnenin voelt \emph{\omterm{def:g4:weight-and-mass:weight}{gewicht}} --- de trek van de appels op de
schaal. Op Aarde reizen de trek en de stof strikt samen (meer trek
precies wanneer meer stof), dus vertelt een weegschaal met een schaal
in kilogrammen stof toch de waarheid --- hier.
\end{solution}

\begin{solution}{exo:g4:weight-and-mass:10}
Metingen van \omterm{def:g3:measuring-mass:mass}{massa} reizen onveranderd mee: ze beschrijven de stof zelf,
niet de plaats. Elke weging die \emph{vergelijkt} --- de
\omterm{def:g3:levers-and-scales:scale}{balansweegschaal} tegen referentieblokjes --- meldt \omterm{def:g3:measuring-mass:mass}{massa} en blijft
bij de verre planeet precies kloppen. Alleen aflezingen van
\omterm{def:g4:weight-and-mass:weight}{gewicht} met een veer krimpen daarginds, en die stonden nooit in de
onderdelenlijst.
\end{solution}

\begin{solution}{exo:g4:weight-and-mass:11}
De \omterm{def:g3:measuring-mass:mass}{massa} is er nog helemaal. Aanwijzingen: de moersleutel heen en
weer schudden kost nog steeds echte moeite --- haar hand voelt de stof
zich tegen het heen en weer verzetten; en weggegooid vliegt hij niet
sneller weg dan haar spieren aankunnen --- een homp stof blijft een
homp stof om te gooien. Het \omterm{def:g4:weight-and-mass:weight}{gewicht} is weg; de stof nooit.
\end{solution}
